\documentclass[12pt]{beamer}
\usepackage[utf8]{inputenc}
\usepackage[T2A]{fontenc}
\usepackage[russian]{babel}
\usepackage{amsmath}
\usepackage{booktabs}
\usepackage{tikz}
\usepackage{pgfplots}
\pgfplotsset{compat=1.18}

\usetheme{Madrid}
\usecolortheme{default}

\title{Мечи залива работорговцев}
\subtitle{Разведывательный анализ данных поставщиков стали}
\author{Студент 1 \and Студент 2 \and Студент 3}
\institute{Практикум, 3 курс}
\date{\today}

\begin{document}

\begin{frame}
    \titlepage
\end{frame}

\begin{frame}{План презентации}
    \tableofcontents
\end{frame}

\section{Цель и задачи}

\begin{frame}{Цель работы}
    \begin{block}{Основная цель}
        На основе данных о производстве и поломках мечей выбрать поставщика стали
        для эксклюзивного контракта: \texttt{harpy.co} или \texttt{westeros.inc}.
    \end{block}
\end{frame}

\begin{frame}{Задачи}
    \begin{enumerate}
        \item Подготовить и очистить данные (обработка аномалий, агрегация по месяцам).
        \item Рассчитать метрики качества и надежности поставщиков.
        \item Учесть тренд изменения брака во времени.
        \item Сформировать итоговую рекомендацию для выбора контракта.
    \end{enumerate}
\end{frame}

\section{Данные и методика}

\begin{frame}{Исходные данные}
    CSV-файл \texttt{production-data.csv} содержит:
    \begin{itemize}
        \item идентификатор кузнеца \texttt{unsullen.id};
        \item месяц производства \texttt{production.date};
        \item месяц отчета \texttt{report.date};
        \item объем выпуска \texttt{produced};
        \item число поломок \texttt{defects};
        \item поставщика стали \texttt{supplier}.
    \end{itemize}

    \vspace{0.3cm}
    Рассматриваются два поставщика: \texttt{harpy.co} и \texttt{westeros.inc}.
\end{frame}

\begin{frame}{Метрики анализа}
    \begin{block}{Базовые метрики}
        \[
        \text{defect\_rate} =
        \frac{\sum \text{defects}}{\sum \text{produced}},
        \quad
        \text{early\_failure\_rate} =
        \frac{\sum \text{defects}_{\;(\text{lag} \le 1)}}{\sum \text{produced}}
        \]
    \end{block}

    \begin{block}{Тренд и стабильность}
        Для месячной доли дефектов строилась линейная модель
        \[
        y_m = a \cdot m + b,
        \]
        где:
        \begin{itemize}
            \item \(a = \text{trend\_slope}\) --- скорость роста/падения брака;
            \item \(\text{residual\_std} = \sigma(y_m - \hat{y}_m)\) --- шум вокруг тренда.
        \end{itemize}
    \end{block}
\end{frame}

\section{Результаты}

\begin{frame}{Сравнение поставщиков по итоговым метрикам}
\small
\begin{center}
\begin{tabular}{lcccccc}
\toprule
Поставщик & defect\_rate & early\_rate & slope & residual\_std & last\_month & worst\_month \\
\midrule
harpy.co      & 0.19282 & 0.02477 & -0.08973 & 0.04052 & 0.02414 & 0.44874 \\
westeros.inc  & 0.26147 & 0.08484 & -0.06769 & 0.01073 & 0.08360 & 0.41421 \\
\bottomrule
\end{tabular}
\end{center}

\vspace{0.2cm}
\footnotesize
\textbf{Интерпретация:}
harpy.co лучше по среднему браку, ранним поломкам, текущему уровню и скорости улучшения;
westeros.inc более стабилен относительно собственного тренда.
\end{frame}

\begin{frame}{Рисунок: доля дефектов (ключевые метрики)}
\centering
\begin{tikzpicture}
\begin{axis}[
    width=0.9\textwidth,
    height=0.55\textwidth,
    ybar,
    bar width=10pt,
    symbolic x coords={harpy,westeros},
    xtick=data,
    ymin=0,
    ylabel={Значение},
    legend style={at={(0.5,-0.2)},anchor=north,legend columns=2},
    nodes near coords,
    nodes near coords align={vertical}
]
\addplot coordinates {(harpy,0.19282) (westeros,0.26147)};
\addplot coordinates {(harpy,0.02477) (westeros,0.08484)};
\legend{defect\_rate,early\_failure\_rate}
\end{axis}
\end{tikzpicture}
\end{frame}

\begin{frame}{Рисунок: тренд месячной доли дефектов}
\centering
\begin{tikzpicture}
\begin{axis}[
    width=0.9\textwidth,
    height=0.55\textwidth,
    xmin=1, xmax=6,
    ymin=0, ymax=0.5,
    xlabel={Месяц производства},
    ylabel={Доля дефектов},
    legend style={at={(0.5,-0.2)},anchor=north,legend columns=2},
    grid=both
]
% Приближенные точки по агрегированным данным из анализа
\addplot+[mark=o, thick] coordinates {
    (1,0.44874) (2,0.34209) (3,0.21629) (4,0.07724) (5,0.04921) (6,0.02414)
};
\addplot+[mark=square, thick] coordinates {
    (1,0.41402) (2,0.36893) (3,0.31276) (4,0.22914) (5,0.15808) (6,0.08360)
};
\legend{harpy.co,westeros.inc}
\end{axis}
\end{tikzpicture}
\end{frame}

\begin{frame}{Рисунок: производство по месяцам}
\centering
\begin{tikzpicture}
\begin{axis}[
    width=0.9\textwidth,
    height=0.55\textwidth,
    xmin=1, xmax=6,
    ymin=5200, ymax=5320,
    xlabel={Месяц производства},
    ylabel={Количество произведенных мечей},
    legend style={at={(0.5,-0.2)},anchor=north,legend columns=2},
    grid=major
]
% Агрегированные значения производства по месяцам
\addplot+[mark=o, thick] coordinates {
    (1,5248) (2,5256) (3,5243) (4,5282) (5,5243) (6,5260)
};
\addplot+[mark=square, thick] coordinates {
    (1,5280) (2,5272) (3,5298) (4,5285) (5,5263) (6,5227)
};
\legend{harpy.co,westeros.inc}
\end{axis}
\end{tikzpicture}
\end{frame}

\section{Выводы}

\begin{frame}{Выводы и рекомендация}
    \begin{enumerate}
        \item У \texttt{harpy.co} ниже средний уровень брака:
        \(0.19282 < 0.26147\).
        \item У \texttt{harpy.co} существенно ниже ранние поломки:
        \(0.02477 < 0.08484\).
        \item Тренд улучшения у \texttt{harpy.co} быстрее (более отрицательный наклон).
        \item Несмотря на более высокий шум вокруг тренда, текущая доля дефектов у
        \texttt{harpy.co} значительно ниже.
    \end{enumerate}

    \vspace{0.3cm}
    \begin{block}{Итог}
        Рекомендуется заключить эксклюзивный контракт с \texttt{harpy.co}.
    \end{block}
\end{frame}

\begin{frame}{Список литературы}
    \begin{thebibliography}{9}
        \bibitem{eda}
        Tukey J.W.\\
        Exploratory Data Analysis. Addison-Wesley, 1977.

        \bibitem{pandas}
        McKinney W.\\
        Python for Data Analysis. O'Reilly, 2022.

        \bibitem{numpy}
        Harris C.R. et al.\\
        Array programming with NumPy. Nature, 2020.

        \bibitem{matplotlib}
        Hunter J.D.\\
        Matplotlib: A 2D Graphics Environment. Computing in Science \& Engineering, 2007.
    \end{thebibliography}
\end{frame}

\end{document}
